\section[Cap\'{i}tulo 5: Herramientas Tecnol\'{o}gicas para el Desarrollo del Prototipo]{\centering Cap\'{i}tulo 5: Herramientas Tecnol\'{o}gicas para el Desarrollo del Prototipo}

Las herramientas tecnol\'{o}gicas son el medio que permite convertir el dise\~{n}o pedag\'{o}gico y l\'{u}dico en un prototipo funcional. Para este proyecto, la elecci\'{o}n de herramientas debe priorizar facilidad de desarrollo, organizaci\'{o}n de escenas, manejo de interfaz y posibilidad de ajustar r\'{a}pidamente los casos de juego (Unity Technologies, 2026).

El prototipo no requiere una tecnolog\'{i}a compleja desde el inicio, sino una base que permita construir pantallas, decisiones, reportes y progresi\'{o}n de niveles. Por ello, se consideran motores y recursos que faciliten iterar la experiencia antes de ampliar su alcance (Godot Engine, 2026).

\subsection{Motores de videojuegos como apoyo}

\begin{center}
\begin{tabular}{|p{0.20\textwidth}|p{0.35\textwidth}|p{0.30\textwidth}|}
\hline
\textbf{Herramienta} & \textbf{Aporte principal} & \textbf{Relaci\'{o}n con el proyecto} \\
\hline
Unity & Organizaci\'{o}n por objetos, escenas, componentes e interfaz (Unity Technologies, 2026). & Adecuado para prototipos con pantallas, botones, reportes y l\'{o}gica de niveles. \\
\hline
Godot & Estructura por nodos y escenas reutilizables (Godot Engine, 2026). & Facilita construir pantallas modulares y casos interactivos ligeros. \\
\hline
GameMaker & Eventos de objetos y flujo visual simple (YoYo Games, 2026). & Puede servir para prototipos 2D centrados en decisiones e interfaz. \\
\hline
\end{tabular}
\end{center}

\subsection{Criterios de selecci\'{o}n}

La herramienta elegida debe permitir que el desarrollo se concentre en la experiencia del jugador y no solamente en resolver problemas t\'{e}cnicos. Para un prototipo acad\'{e}mico, la facilidad para crear interfaces, registrar decisiones y modificar casos es m\'{a}s importante que usar funciones avanzadas del motor (Godot Engine, 2026).

Tambi\'{e}n se debe considerar la disponibilidad de documentaci\'{o}n, compatibilidad con recursos visuales y facilidad para exportar una versi\'{o}n demostrable. Estos criterios ayudan a que el prototipo pueda probarse con usuarios y ajustarse seg\'{u}n observaciones de usabilidad (Unity Technologies, 2026).

En el caso del videojuego propuesto, las herramientas deben apoyar una estructura modular: men\'{u}s, niveles, casos, registro de decisiones y pantalla final de reporte. Esa organizaci\'{o}n permite ampliar o corregir escenarios sin rehacer todo el sistema (YoYo Games, 2026).

\subsection{S\'{i}ntesis del cap\'{i}tulo}

Las herramientas tecnol\'{o}gicas no son el centro del proyecto, pero s\'{i} condicionan la forma en que se construye el prototipo. Una buena selecci\'{o}n debe permitir desarrollar r\'{a}pido, probar escenarios y mantener una estructura clara para futuras mejoras.
